% =====================================================================
%  Section 10 -- Discussion (owner: LEAD).
% =====================================================================
\section{Discussion}\label{sec:discussion}

\paragraph{What a well-designed audit could establish}
Debate about recommender systems and hateful content tends to ask whether
``the algorithm'' amplifies it. The framework replaces that question with
several that have precise answers under stated assumptions. For a stated
service, recommendation surface, jurisdiction and audit window, and for each
confirmatory persona, an audit run under the protocol of \cref{sec:protocol}
would report the amplification ratio and the additive amplification relative
to a named reference policy (\cref{def:amplification}), corrected for
classifier error in each arm (\cref{thm:misclass}), with a confidence sequence
that is approximately valid at whatever point the audit stopped
(\cref{alg:sequential,thm:sequential}), and a conclusion fixed in advance:
material amplification, material de-amplification, no material difference at
the registered margin, a difference of undetermined materiality, or an
inconclusive result (\cref{tab:prot-decisions}). Only the four-persona
materiality audit, rung R1 of the feasibility gate, would also report the
persona-contrast decomposition (\cref{sec:est-audit}).
% ALIGN-WITH-SIM: status of the Markov summaries, aligned with E4b (07, sec:sim-markov).
The Markov summaries of each arm (\cref{prop:markov,thm:perturb}) would be
reported as exploratory and model-based. Their value is that they reach beyond
the audited session length, where short sessions can understate the effect of
a change many times over (\cref{sec:markov-perturb}), but the protocol's slate
states are observed through a classifier: in simulation the uncorrected chain
can reverse the comparison of the arms, and only the corrected two-state chain
was reliable (\cref{sec:sim-markov}).

What the audit would in fact establish is smaller, and the paper's power
calculations say so. Under plausible inputs, the four-persona materiality
audit needs about $4{,}900$ accounts and $170{,}000$ coded exposures, and
under illustrative capacity caps the feasibility gate selects rung R4: the
detection of amplification, or de-amplification, for the Jewish-interest
persona alone, with ``difference detected; materiality undetermined'' as the
conclusion it is powered to reach. That is a modest result, but a defined
one, and the alternative, an audit too small to answer its question, would
spend its footprint on no knowledge (\cref{sec:prot-power}).

Such an audit could not establish everything that is asked of it, and the
framework says exactly what it leaves out. A fixed-demand audit says nothing
about how the algorithm shapes what people want, and without outside
information on how demand forms, the total effect of the algorithm on real
users is not identified (\cref{prop:audit-scope}). A finding that
$\mathrm{AR}(d)>1$ for one persona is a statement about supply at that
persona's demand. It is not a measurement of how much antisemitic content
real users see, nor evidence about other surfaces, other services or other
periods, nor about the intentions of the people who built the system or the
legality of its conduct (\cref{sec:eth-neutrality}). A finding that
$\mathrm{AR}(d)\le1$ is limited in the same ways. Within those limits, an
estimate of this kind is what a systemic-risk assessment under Article~34 of
the DSA could be checked against: a provider's statement that its recommender
does not raise exposure to antisemitic content is a claim about supply at
given demand, and it is the claim an independent audit can test
(\cref{sec:bg-governance}).

\paragraph{Limitations}
We stress six.

\emph{The moving target.} An audit estimates amplification by the policies
that ran during its window, and the estimate expires with the system it
describes. Blocking keeps drift common to both arms out of the contrast, and
sentinels and a test across waves make drift visible (\cref{sec:id-drift}),
but nothing makes a result from one window a result about the next, as the
period qualifier of \citet{hosseinmardi2024} illustrates. A service that
learns to recognise audit accounts breaks transport from puppets to people
(\cref{asm:id-transport}), and the audit cannot test for it.

\emph{From puppets to people.} A persona is a script; real users have
histories, social graphs and preferences that respond to what they are shown.
The step to real users rests on \cref{asm:id-transport}, which audit data
cannot test. The choice rule adds a gap by design: puppets never open a
flagged item, so $\mathrm{AR}(d)$ describes the supply to a demand profile
that does not engage with antisemitic content, the right estimand for users
who do not seek such content, not for those who do. Population statements
need weights over demand profiles, which an audit does not supply.

\emph{Dependence on the classifier and the codebook.} The true label is a
coding decision that depends on the codebook and the definition of
antisemitism behind it, which is why scope variants are reported side by
side. The correction removes classifier bias only if the error rates hold
within the strata where they are estimated, which emerging coded language
\citep{zannettou2020}, quotation and counter-speech \citep{rottger2021} and
images and video can each break. Classifier error is modelled here, not
eliminated.

\emph{Law and access.} Terms of service, contract law and
computer-crime statutes constrain sock-puppet audits (\cref{sec:eth-legal}),
and the practical reach of vetted access under Article~40(4) is still being
established, so the services that are audited are not a random sample of
services.

\emph{Heterogeneity across services.} Surfaces and services differ in their
reference policies: by
\cref{lem:est-agg}(c), changing the reference multiplies the ratio by a
factor that does not involve the algorithm, so ratios are comparable across
services only on a common reference or a stated ladder of references, with
the same classifier, codebook and personas. Interference differs too, and
because $\xi$ is not identified, each conclusion carries its breakdown value
rather than a corrected estimate (\cref{prop:interference}).

\emph{Feasibility.} The audit the questions call for is not the audit a
small lab can run. A materiality test across four personas needs thousands of
accounts and coded exposures (\cref{tab:prot-power}), and the protocol
therefore commits only to a pilot and a registered gate, which under
illustrative caps selects detection for one persona. The decomposition across
personas, materiality conclusions and evidence of no material difference are
then out of reach. Platform-side access under Article~40(4), where it can
deliver exposure under a fixed demand profile, removes the need for accounts
and with it most of this constraint; until it does, single-persona detection
audits are what independent researchers can responsibly run. The
approximations add to this: the sequential guarantee is asymptotic, and at
the registered sizes its level is checked by simulation and by a pilot whose
ten puppets per cell cannot reveal the tails of the double-sampled outcome
(\cref{sec:prot-stopping}).
% ALIGN-WITH-SIM: E5b finding on the primary sequence, aligned (07, sec:sim-sequential).
In simulation the sequence on $\widetilde E_u$ was conservative at the
registered design despite heavy tails, but that is evidence about one
synthetic world, not a guarantee for a real audit.

\paragraph{Extensions}
Three extensions seem most valuable. For \emph{generative answer engines},
the policy is the system that generates the answer, a persona is a script of
prompts, exposure is the share of answers that a validated classifier labels
antisemitic, and a natural reference is the same system without generation;
generation is stochastic at fixed inputs and has no fixed eligible pool, so
the distinction between content and algorithmic supply must be redrawn. In
\emph{search}, demand is explicit in the query, a persona becomes a
distribution over queries, and the same query without a profile is a natural
reference \citep{robertson2023}. \emph{Cross-platform audits} with harmonised
personas, one frozen classifier and a shared validation sample would make
estimates comparable across services, and adaptive personas would let them
study demand shaping for synthetic users (\cref{sec:est-audit}).

\paragraph{Why define amplification}
A count of antisemitic items seen online says how much there is. It does not
say where it came from. Content that users seek out calls for one kind of
response, content that a recommender supplies beyond what users seek calls
for another, and content that a recommender has taught users to seek calls
for a third. These are different claims about responsibility, and they carry
different duties under the law and different obligations for the people who
design these systems. An audit can test the second claim, for a stated system
and a stated alternative, and it can say plainly that it cannot test the
third. Stating the comparison before measuring it, correcting for the errors
of the instrument, stopping as soon as the evidence allows and fixing the
rules before the data are seen are how an audit can take seriously both the
people who are targeted by antisemitic content and the evidence about where
it comes from.
